\thispagestyle{empty}
\begin{tikzpicture}[remember picture,overlay]
\begin{scope}[shift={(current page.south west)},x=1mm,y=1mm]
\fill[otNavy](0,0)rectangle(210,297);
\fill[otNavyDark](0,0)rectangle(210,88);
\draw[white,opacity=.055,step=7mm](0,0)grid(210,297);
\node[anchor=north west,inner sep=0pt]at(23,278){\includegraphics[width=57mm]{images/logo-blanc.png}};
\node[anchor=west,text=otGreen,font=\otSemi\fontsize{12}{15}\selectfont]at(24,235){MPSI\quad /\quad PHYSIQUE};
\node[anchor=west,text=white,font=\sffamily\bfseries\fontsize{52}{56}\selectfont]at(22,211){Physique};
\node[anchor=north west,text=white,text width=158mm,font=\otLight\fontsize{31}{36}\selectfont]at(24,191){Formation des images\\et optique géométrique};
\fill[otGreen](24,151)rectangle(54,152.4);
\node[anchor=north west,text=white,text width=153mm,font=\sffamily\fontsize{11}{16}\selectfont]at(24,139){Cours développé et schémas vectoriels\par Exercices progressifs et corrigés expliqués\par Défis avancés, TP, DM et DS avec barèmes\par Parcours sur des sujets authentiques};
\draw[white,opacity=.65,line width=.8pt](22,50)--(188,50);
\draw[white,opacity=.7,line width=1.2pt,{Stealth}-{Stealth}](98,14)--(98,84);
\draw[otGreen,line width=1.7pt](22,74)--(98,74)--(160,24);
\draw[white,opacity=.7,line width=1pt](22,24)--(98,50)--(188,80.8);
\node[anchor=west,text=white,font=\otSemi\fontsize{9}{12}\selectfont]at(24,8){CHAPITRE 01\quad /\quad VERSION À VALIDER\quad /\quad 25 OCTOBRE 2026};
\end{scope}
\end{tikzpicture}
\null

\Page{Collection OptiTAB}{Physique MPSI}{page-titre}
\thispagestyle{empty}
\vspace{9mm}
\begin{center}
\includegraphics[width=63mm]{images/logo-couleur.png}\par
\vspace{19mm}
{\otSemi\fontsize{11}{15}\selectfont\color{otGreenDark}CLASSE PRÉPARATOIRE SCIENTIFIQUE}\par
\vspace{5mm}
{\otLight\fontsize{23}{29}\selectfont Cours et exercices corrigés}\par
\vspace{10mm}
{\color{otGreen}\rule{25mm}{1.2pt}}\par
\vspace{9mm}
{\fontsize{15}{21}\selectfont Formation des images\par et optique géométrique}\par
\vspace{8mm}
{\fontsize{11}{16}\selectfont 24 exercices progressifs\par 3 défis avancés et leurs corrigés\par TP, activité numérique, DM et DS barémés}\par
\vspace{13mm}
{\sffamily\small\color{otMuted}Chapitre 01\quad /\quad Édition pilote à valider}\par
\vspace{6mm}
{\fontsize{11}{15}\selectfont Collection OptiTAB}\par
{\sffamily\small\color{otMuted}\href{https://www.optitab.net}{optitab.net}}
\end{center}

\Page{Édition et références}{Mentions éditoriales}{mentions}
\thispagestyle{mpsiFront}
\vspace{9mm}
\begin{minipage}{137mm}
\textbf{Physique MPSI}\par
Formation des images et optique géométrique\par
Chapitre 01 du projet de manuel annuel\par
\vspace{3mm}
Conception des contenus et mise en page : collection OptiTAB.\par
Charte graphique commune aux livres OptiTAB de physique-chimie : typographie, couleurs, repères de méthode et navigation entre énoncés et corrigés.\par
\vspace{3mm}
Édition pilote du 25 octobre 2026, révisée pour compléter les pages liminaires. Le cours et les entraînements du premier chapitre sont disponibles ; les autres chapitres figurent dans le plan de l'ouvrage pour préparer la suite.\par
\vspace{3mm}
\textbf{Référence pédagogique}\par
Programme de physique-chimie de la voie MPSI, annexe 2 de l'arrêté du 5 janvier 2021, NOR ESRS2035779A, publié au Bulletin officiel spécial n° 1 du 11 février 2021. Ce volume est consacré aux thèmes de physique et à la formation expérimentale de ce programme.\par
Le premier chapitre traite la section 1.1 « Formation des images ». La concordance des capacités est donnée p.\,\pageref{concordance}. L'ordre et le découpage des chapitres constituent des choix pédagogiques de la collection.\par
\vspace{3mm}
\textbf{Énoncés, figures et données}\par
Les exercices et évaluations portant la mention OptiTAB sont originaux. Les schémas sont construits dans le modèle indiqué ; une représentation schématique ou sans échelle est signalée lorsque cela est utile. Les séries créées pour l'entraînement ne sont pas présentées comme des mesures réellement réalisées.\par
Les supports authentiques extérieurs sont identifiés par leur diffuseur, leur date, la partie travaillée et leur lien d'accès. Les références et conditions de diffusion restent celles des auteurs ou organismes concernés.\par
\vspace{3mm}
© OptiTAB, 2026.\par
\href{https://www.optitab.net}{optitab.net}
\end{minipage}

\Page{Pour ouvrir le livre}{Avant-propos}{avant-propos}
\thispagestyle{mpsiFront}
\begin{minipage}{140mm}
{\fontsize{14}{20}\selectfont\itshape\color{otNavy}
Ce livre accompagne le passage du lycée à la physique de MPSI : comprendre les modèles, raisonner avec précision et apprendre à chercher de manière autonome.\par}
\vspace{4mm}
La première année de classe préparatoire demande de relier des connaissances qui pouvaient sembler séparées : une observation, un schéma, une équation, une mesure et une simulation décrivent souvent le même phénomène sous des formes différentes. L'objectif est d'apprendre à passer de l'une à l'autre en gardant le sens physique.\par
\vspace{2mm}
Chaque chapitre suit une progression commune. Un diagnostic réactive les acquis ; le cours développe les lois et leurs conditions d'application ; les méthodes montrent comment conduire une résolution. Les schémas soutiennent le raisonnement et permettent de vérifier les signes, les directions et les ordres de grandeur.\par
\vspace{2mm}
Les exercices vont des applications directes aux problèmes demandant davantage d'initiative. Les défis proposent des situations plus ouvertes, accompagnées d'indices progressifs. Les travaux pratiques et les activités numériques entraînent à confronter un modèle à des résultats et à discuter la précision d'une conclusion.\par
\vspace{2mm}
Les corrigés expliquent les choix de méthode, les calculs et l'interprétation des résultats. Ils servent à comprendre une difficulté après une recherche personnelle. Refaire un exercice quelques jours plus tard, sans regarder sa solution, reste une étape essentielle du travail.\par
\vspace{2mm}
Les DM et les DS permettent de travailler la rédaction et la gestion du temps. Des parcours sur de vrais devoirs et sujets de concours complètent l'entraînement ; les références permettent de retrouver le support d'origine et de distinguer ces documents des créations OptiTAB.\par
\vspace{2mm}
Cette édition pilote présente le premier chapitre complet et le plan du futur ouvrage. Elle permet de valider ensemble le niveau d'explication, la place des figures, la progression des exercices et la présentation avant de poursuivre. Le tableau de bord t'aide dès maintenant à repérer ce qui est acquis et ce qu'il faut reprendre.\par
\vspace{5mm}
{\sffamily\small\color{otNavy}Bon travail et bonne progression en physique.}\par
\vspace{2mm}
{\sffamily\small\color{otMuted}La collection OptiTAB}
\end{minipage}


\Page{Pour commencer}{Utiliser ce chapitre}{mode-emploi}
\thispagestyle{mpsiFront}
Ce chapitre est le premier jalon d'un manuel de \textbf{physique MPSI}. Il développe le bloc « Formation des images » du programme national. La suite du livre sera rédigée après validation de ce pilote.

\begin{loi}{Le contrat de travail}
Comprendre les hypothèses, construire le schéma, choisir une loi, conduire le calcul puis discuter le résultat. Chaque solution doit pouvoir être expliquée à l'oral, sans se réduire à une formule apprise.
\end{loi}
\Sub{Un parcours en quatre étapes}
\begin{enumerate}
\item Faire le diagnostic p.\,\pageref{diagnostic}, puis lire le cours p.\,\pageref{c-sources} à \pageref{c-synthese}. Refaire les démonstrations signalées.
\item Résoudre les 24 exercices p.\,\pageref{td1} à \pageref{td8}. Les durées sont indicatives ; la difficulté augmente de 1 à 3.
\item Réaliser le TP, puis traiter les trois défis, le DM et le DS. Les indices p.\,\pageref{indices} permettent de repartir sans lire immédiatement la solution.
\item Corriger en explicitant l'erreur commise. Les solutions commencent p.\,\pageref{corr-diag}. Les parcours authentiques sont p.\,\pageref{authentiques}.
\end{enumerate}
\Sub{Repères rapides}
\begin{tabularx}{\linewidth}{@{}Y r@{}}
\toprule
Programme et concordance & \pageref{concordance}\\
Cours, méthodes et schémas & \pageref{c-sources}\\
Expérience et calcul numérique & \pageref{c-tp}\\
Exercices et problèmes & \pageref{td1}\\
Défis et évaluations & \pageref{defis1}\\
Corrigés et barèmes & \pageref{corr-diag}\\
Supports réels et bibliographie & \pageref{authentiques}\\
\bottomrule
\end{tabularx}
\source{Convention éditoriale : les énoncés portant la mention « OptiTAB » sont originaux. Les supports authentiques sont identifiés par auteur ou diffuseur, date, partie et lien. Les corrigés de ce chapitre sont rédigés indépendamment.}
\source{Référence réglementaire : \href{https://www.education.gouv.fr/bo/21/Special1/ESRS2035779A.htm}{arrêté du 5 janvier 2021, NOR ESRS2035779A}, BO spécial n° 1 du 11 février 2021, annexe 2. Recherche documentaire vérifiée le 25 octobre 2026.}

\Page{Organiser son travail}{Parcours de travail}{parcours}
\thispagestyle{mpsiFront}
La progression du livre peut accompagner celle du professeur. Les outils de mesure, les incertitudes, le calcul et la programmation se travaillent dès les premières séances, au fil des thèmes.
\Sub{Choisir son point de départ}
\begingroup\renewcommand{\arraystretch}{1.28}
\begin{tabularx}{\linewidth}{@{}p{35mm}Y@{}}
\toprule
Parcours & Travail proposé dans le premier chapitre\\\midrule
Reprendre les bases & Diagnostic, cours, puis E01 à E09. Refaire les schémas et annoncer les hypothèses.\\[2mm]
Installer les méthodes & E10 à E18 : conjugaison, signes, foyers, projection, œil et appareil photographique.\\[2mm]
Gagner en autonomie & E19 à E24 : réglages, fibre, association de lentilles et incertitudes.\\[2mm]
Approfondir & Défis D1 à D3. Chercher d'abord seul ; consulter les indices progressivement.\\[2mm]
Préparer les devoirs & DM, puis DS de 2 h. Corriger avec le barème et reprendre les capacités fragiles.\\[2mm]
Se confronter à un support réel & Parcours R1 à R3 : retrouver le document original et traiter les questions ciblées.\\
\bottomrule
\end{tabularx}\endgroup
\Sub{Une séance de travail en quatre temps}
\begin{enumerate}
\item \textbf{Reformuler.} Dire avec ses propres mots ce qu'énonce une loi, avec ses conditions de validité et les unités des grandeurs.
\item \textbf{Chercher.} Dessiner la situation, nommer les inconnues et choisir une stratégie avant de remplacer les lettres par des nombres.
\item \textbf{Corriger.} Comparer les étapes du raisonnement ; noter la première erreur ou la première idée manquante.
\item \textbf{Réactiver.} Refaire la question quelques jours plus tard et expliquer la solution à voix haute, comme en interrogation orale.
\end{enumerate}
\begin{methode}{Un rythme à adapter}
Pour ce premier chapitre : consacrer une première étape au cours et aux constructions, une deuxième aux exercices et au TP, puis une troisième aux défis et évaluations. La compréhension et la capacité à refaire seul déterminent le passage à l'étape suivante.
\end{methode}
\source{Les difficultés et durées indiquées dans les exercices servent de repères. Le barème d'un DS est un outil pédagogique de ce livre, pas un barème officiel de concours.}

\Page{Collection OptiTAB}{Sommaire}{sommaire-general}
\thispagestyle{mpsiFront}
\begingroup\setlength{\parskip}{0pt}\fontsize{10}{13.5}\selectfont
{\sffamily\fontsize{9}{12}\selectfont\color{otMuted}MPSI\enspace·\enspace Physique\enspace·\enspace Plan de l'ouvrage}\par\vspace{3mm}
\tocpart{1}{Ondes et signaux}
\planitem{01}{Formation des images et optique géométrique}{1.1}
\planitem{02}{Signaux électriques dans l'ARQS}{1.2}
\planitem{03}{Circuits linéaires du premier ordre}{1.3}
\planitem{04}{L'oscillateur harmonique}{1.4}
\planitem{05}{Oscillations libres amorties}{1.4}
\planitem{06}{Régime sinusoïdal forcé et résonance}{1.4}
\planitem{07}{Filtrage linéaire}{1.5}
\planitem{08}{Propagation d'un signal}{1.6}
\planitem{09}{Superposition et interférences}{1.6}
\planitem{10}{Ondes stationnaires et corde vibrante}{1.6}
\tocpart{2}{Mouvements et interactions}
\planitem{11}{Décrire le mouvement d'un point}{2.1}
\planitem{12}{Lois de Newton et dynamique}{2.2}
\planitem{13}{Énergie du point matériel}{2.3}
\planitem{14}{Particules chargées dans des champs uniformes}{2.4}
\planitem{15}{Moment cinétique}{2.5}
\planitem{16}{Forces centrales et satellites}{2.6}
\planitem{17}{Mouvement d'un solide}{2.7}
\tocpart{3}{L'énergie : conversions et transferts}
\planitem{18}{Décrire un système à l'équilibre}{3.1}
\planitem{19}{Transformations, travail et transfert thermique}{3.2}
\planitem{20}{Premier principe et bilans d'énergie}{3.3}
\planitem{21}{Deuxième principe et bilans d'entropie}{3.4}
\planitem{22}{Machines thermiques}{3.5}
\tocpart{4}{Induction et physique quantique}
\planitem{23}{Champ magnétique}{1.7.1}
\planitem{24}{Actions d'un champ magnétique}{1.7.2}
\planitem{25}{Induction dans un circuit fixe}{1.7.3 et 1.7.4}
\planitem{26}{Induction et circuits mobiles}{1.7.5}
\planitem{27}{Introduction à la physique quantique}{1.8}
\tocpart{5}{Outils transversaux}
\planitem{28}{Mesures, incertitudes et méthodes expérimentales}{Formation expérimentale}

\par\vspace{2mm}{\color{otRule}\rule{\linewidth}{.4pt}}\par\vspace{2mm}
\tocresource{Défis, DM et DS}{defis1}
\tocresource{Corrigés}{corr-diag}
\tocresource{Supports authentiques et références}{authentiques}
\vspace{1mm}
\source{Chapitre 01 disponible\enspace·\enspace « - » : chapitre en préparation.}
\endgroup

\Page{Collection OptiTAB}{Parcours et ressources}{sommaire-ressources}
\thispagestyle{mpsiFront}
\source{Compléments, entraînements et outils de révision du chapitre 01. Les liens renvoient aux pages présentes dans cette édition.}
\Sub{Approfondir}
\tocresource{Défi D1\quad Profondeur de champ}{defis1}
\tocresource{Défi D2\quad Liaison optique}{defis2}
\tocresource{Défi D3\quad Microscope}{defis3}
\tocresource{Indications progressives}{indices}
\Sub{Expérimenter et s'évaluer}
\tocresource{TP\quad Former et mesurer une image}{c-tp}
\tocresource{Activité numérique\quad Incertitude et simulation}{c-numerique}
\tocresource{Devoir maison\quad Focométrie de Bessel}{dm}
\tocresource{Devoir surveillé\quad Sujet de 2 heures}{ds1}
\Sub{Comprendre et corriger}
\tocresource{Corrigé du diagnostic}{corr-diag}
\tocresource{Corrigés des 24 exercices}{corr-e01}
\tocresource{Corrigés des trois défis}{corr-d1}
\tocresource{Corrigé et barème du DM}{corr-dm}
\tocresource{Corrigé et barème du DS}{corr-ds1}
\Sub{Se repérer et travailler sur des sujets réels}
\tocresource{Tableau de progression}{tableau-de-bord}
\tocresource{Notions et capacités du programme}{concordance}
\tocresource{Parcours R1 et R2\quad DM et DS de MPSI}{authentiques}
\tocresource{Parcours R3\quad Mines-Ponts MP 2021}{auth-concours}
\tocresource{Références et suivi}{references}

\Page{Suivre sa progression}{Mon tableau de bord}{tableau-de-bord}
\thispagestyle{mpsiFront}
Après chaque exercice, cocher la situation qui correspond à la dernière tentative : \textbf{S} = réussi seul ; \textbf{A} = réussi avec une aide ; \textbf{R} = à refaire. Une nouvelle tentative peut faire évoluer ce bilan.\par
\vspace{2mm}
\begin{minipage}[t]{.47\linewidth}
\renewcommand{\arraystretch}{1.2}
\begin{tabularx}{\linewidth}{@{}Yccc@{}}\toprule
Exercice&S&A&R\\\midrule
\hyperref[ex-E01]{E01}&\caseSuivi&\caseSuivi&\caseSuivi\\
\hyperref[ex-E02]{E02}&\caseSuivi&\caseSuivi&\caseSuivi\\
\hyperref[ex-E03]{E03}&\caseSuivi&\caseSuivi&\caseSuivi\\
\hyperref[ex-E04]{E04}&\caseSuivi&\caseSuivi&\caseSuivi\\
\hyperref[ex-E05]{E05}&\caseSuivi&\caseSuivi&\caseSuivi\\
\hyperref[ex-E06]{E06}&\caseSuivi&\caseSuivi&\caseSuivi\\
\hyperref[ex-E07]{E07}&\caseSuivi&\caseSuivi&\caseSuivi\\
\hyperref[ex-E08]{E08}&\caseSuivi&\caseSuivi&\caseSuivi\\
\hyperref[ex-E09]{E09}&\caseSuivi&\caseSuivi&\caseSuivi\\
\hyperref[ex-E10]{E10}&\caseSuivi&\caseSuivi&\caseSuivi\\
\hyperref[ex-E11]{E11}&\caseSuivi&\caseSuivi&\caseSuivi\\
\hyperref[ex-E12]{E12}&\caseSuivi&\caseSuivi&\caseSuivi\\\bottomrule
\end{tabularx}
\end{minipage}\hfill
\begin{minipage}[t]{.47\linewidth}
\renewcommand{\arraystretch}{1.2}
\begin{tabularx}{\linewidth}{@{}Yccc@{}}\toprule
Exercice&S&A&R\\\midrule
\hyperref[ex-E13]{E13}&\caseSuivi&\caseSuivi&\caseSuivi\\
\hyperref[ex-E14]{E14}&\caseSuivi&\caseSuivi&\caseSuivi\\
\hyperref[ex-E15]{E15}&\caseSuivi&\caseSuivi&\caseSuivi\\
\hyperref[ex-E16]{E16}&\caseSuivi&\caseSuivi&\caseSuivi\\
\hyperref[ex-E17]{E17}&\caseSuivi&\caseSuivi&\caseSuivi\\
\hyperref[ex-E18]{E18}&\caseSuivi&\caseSuivi&\caseSuivi\\
\hyperref[ex-E19]{E19}&\caseSuivi&\caseSuivi&\caseSuivi\\
\hyperref[ex-E20]{E20}&\caseSuivi&\caseSuivi&\caseSuivi\\
\hyperref[ex-E21]{E21}&\caseSuivi&\caseSuivi&\caseSuivi\\
\hyperref[ex-E22]{E22}&\caseSuivi&\caseSuivi&\caseSuivi\\
\hyperref[ex-E23]{E23}&\caseSuivi&\caseSuivi&\caseSuivi\\
\hyperref[ex-E24]{E24}&\caseSuivi&\caseSuivi&\caseSuivi\\\bottomrule
\end{tabularx}
\end{minipage}
\Sub{Approfondissements et évaluations}
\renewcommand{\arraystretch}{1.2}
\begin{tabularx}{\linewidth}{@{}Yccc@{}}\toprule
Travail&S&A&R\\\midrule
\hyperref[defis1]{Défi D1\quad Profondeur de champ}&\caseSuivi&\caseSuivi&\caseSuivi\\
\hyperref[defis2]{Défi D2\quad Liaison optique}&\caseSuivi&\caseSuivi&\caseSuivi\\
\hyperref[defis3]{Défi D3\quad Microscope}&\caseSuivi&\caseSuivi&\caseSuivi\\
\hyperref[dm]{DM\quad Focométrie de Bessel}&\caseSuivi&\caseSuivi&\caseSuivi\\
\hyperref[ds1]{DS\quad Bilan du chapitre}&\caseSuivi&\caseSuivi&\caseSuivi\\\bottomrule
\end{tabularx}
\Sub{Ma prochaine séance}
Date : \rule{30mm}{.3pt}\qquad Exercice ou capacité à reprendre :\par
\vspace{2mm}\rule{\linewidth}{.3pt}\par
L'erreur ou l'idée à retenir :\par
\vspace{2mm}\rule{\linewidth}{.3pt}

\Page{Chapitre 01\quad /\quad Ondes et signaux}{Formation des images}{chapitre-01}
\vspace{2mm}
{\otLight\fontsize{18}{23}\selectfont\color{otNavy}Optique géométrique, instruments et guidage}\par
\vspace{3mm}
Comment une lentille peut-elle former une image sur un écran, donner une image virtuelle ou guider la lumière vers un capteur ? Le même modèle de rayons permet d'étudier ces situations, à condition de connaître ses hypothèses et ses limites.
\Sub{Les objectifs du chapitre}
\begin{itemize}
\item Caractériser une source lumineuse et relier fréquence, longueur d'onde, vitesse et indice.
\item Appliquer les lois de Snell-Descartes, établir un seuil de réflexion totale et tracer le trajet d'un rayon.
\item Construire et calculer des images par un miroir plan ou une lentille mince, avec des distances algébriques cohérentes.
\item Relier la netteté à un détecteur et comprendre les réglages d'un œil ou d'un appareil photographique.
\item Établir les conditions de guidage et l'étalement temporel d'une fibre à saut d'indice.
\item Former une image au laboratoire, exploiter des mesures et discuter une incertitude.
\end{itemize}
\Sub{Les acquis à réactiver}
Conversions d'unités, trigonométrie du triangle rectangle, angles en radians, proportionnalité et résolution d'équations simples. Le diagnostic p.\,\pageref{diagnostic} permet de faire le point avant de commencer.
\begin{loi}{Le parcours du chapitre}
\begin{tabularx}{\linewidth}{@{}Y r@{}}
\tocrow{Notions et capacités du programme}{concordance}
\tocrow{Cours et schémas}{c-sources}
\tocrow{TP et activité numérique}{c-tp}
\tocrow{24 exercices progressifs}{td1}
\tocrow{3 défis, DM et DS}{defis1}
\tocrow{Indications puis corrigés}{indices}
\tocrow{Parcours sur des sujets authentiques}{authentiques}
\end{tabularx}
\end{loi}


\Page{Programme officiel}{Ce que tu dois savoir faire}{concordance}
La concordance ci-dessous couvre la section \textbf{1.1 Formation des images}, pages internes 12 et 13 de l'annexe de physique-chimie MPSI. Le découpage en chapitres est un choix éditorial ; il n'est pas imposé par le texte officiel.
\begingroup\small\renewcommand{\arraystretch}{1.35}
\begin{tabularx}{\linewidth}{@{}Y p{43mm}@{}}
\toprule
Capacité travaillée & Où la vérifier\\\midrule
Décrire un spectre et relier couleur et longueur d'onde & Cours p.\,\pageref{c-sources} ; E01, E02\\
Définir le modèle des rayons et ses limites & p.\,\pageref{c-modele} ; E03\\
Appliquer Snell-Descartes et établir la réflexion totale & p.\,\pageref{c-lois}--\pageref{c-totale} ; E04--E06\\
Construire une image par un miroir plan & p.\,\pageref{c-miroir} ; E07\\
Énoncer les conditions de Gauss et relier netteté et détecteur & p.\,\pageref{c-gauss} ; E08, E20\\
Définir centre, foyers principaux et secondaires, focale, vergence & p.\,\pageref{c-rayons} ; E09, E13\\
Construire des images réelles ou virtuelles, à distance finie ou infinie & p.\,\pageref{c-reelle}--\pageref{c-virtuelle} ; E10--E13\\
Employer Descartes et Newton, calculer le grandissement & p.\,\pageref{c-newton} ; E10, E14\\
Établir la condition de projection d'un objet réel & p.\,\pageref{c-projection} ; E15, DM\\
Modéliser l'œil, citer résolution et accommodation & p.\,\pageref{c-oeil} ; E16, E17\\
Modéliser un appareil photographique et ses réglages & p.\,\pageref{c-photo} ; E18, E19\\
Construire géométriquement la profondeur de champ & p.\,\pageref{c-profondeur} ; E20, D1\\
Établir cône d'acceptance et dispersion intermodale & p.\,\pageref{c-fibre}--\pageref{c-dispersion} ; E21, E22, D2\\
Former et mesurer une image ; exploiter les incertitudes & TP et numérique p.\,\pageref{c-tp}--\pageref{c-numerique}\\
\bottomrule
\end{tabularx}\endgroup
\begin{methode}{Approfondissements identifiés}
La méthode de Bessel, les associations de lentilles et le microscope sont des applications des outils du programme. Les défis et les parcours de concours sont guidés ; leurs résultats spécifiques ne sont pas des formules supplémentaires à mémoriser.
\end{methode}

\Page{Diagnostic\quad /\quad 15 min}{Avant d'ouvrir le cours}{diagnostic}
Sans document, répondre en une ou deux phrases ou par un calcul court. Ce diagnostic n'est pas une note de niveau : il identifie les acquis à réactiver.
\begin{enumerate}[label=\textcolor{otNavy}{\bfseries D\arabic*.},itemsep=5mm]
\item Une lumière de fréquence $5{,}00\times10^{14}\ \mathrm{Hz}$ se propage dans le vide. Calculer sa longueur d'onde avec $c=3{,}00\times10^8\ \mathrm{m\,s^{-1}}$.
\item Un rayon atteint une surface sous un angle de $30^\circ$ avec la surface. Quel angle faut-il introduire dans la loi de réfraction ?
\item La fréquence d'une onde lumineuse change-t-elle quand elle traverse un dioptre immobile ? Et sa longueur d'onde ?
\item Une lentille convergente de focale $10\ \mathrm{cm}$ a quelle vergence en dioptries ?
\item Un objet est à $30\ \mathrm{cm}$ à gauche d'une lentille. L'axe est orienté vers la droite. Donner la valeur algébrique de $\oa$.
\item Une image virtuelle peut-elle être photographiée ? Peut-on la recueillir directement sur un écran placé à sa position ? Expliquer la différence.
\item Pour un petit angle de $2^\circ$, calculer sa valeur en radians puis estimer sa tangente.
\item Quel élément de l'œil joue le rôle de capteur ? Quel paramètre varie lors de l'accommodation dans le modèle de l'œil réduit ?
\end{enumerate}
\begin{docu}{Bilan personnel}
\checkline{Les conversions et les angles sont maîtrisés.}
\checkline{Je distingue longueur positive et distance algébrique.}
\checkline{Je sais expliquer ce qu'est une image virtuelle.}
\checkline{Je peux dessiner une lentille et ses foyers.}
Corrigé p.\,\pageref{corr-diag}. Une hésitation sur D2 ou D5 mérite d'être corrigée avant les exercices.
\end{docu}
